\clearpage
\onecolumn
\thispagestyle{empty}
\vspace*{\fill}
\begin{center}
	{\color{radblue}\rule{\linewidth}{2pt}}\par
	\vspace{6mm}
	{\Huge\bfseries\color{radblue} Vías Aéreas Superiores}\par
	\vspace{6mm}
	{\color{radblue}\rule{\linewidth}{2pt}}
\end{center}
\vspace*{\fill}
\clearpage
\twocolumn
\pagestyle{contenido}



	
\clearpage	
	%%%%%%%%%%%%%%%%%%%%%%%%%%%%%%%%%%%%%%%%%%%%%%%%%%%%%%%%%%%
	%  SECCIÓN: Vías Aéreas
	%  Archivo: secciones/vias_aereas.tex
	%  Sin preámbulo ni \begin{document} — lo maneja main.tex
		%%%%%%%%%%%%%%%%%%%%%%%%%%%%%%%%%%%%%%%%%%%%%%%%%%%%%%%%%%%
		\onecolumn
		
		\section{Vías Aéreas}
		
		\subsection{Vías Respiratorias Superiores}
		
		\subsubsection*{Laringe y Tráquea}
		
		\paragraph{Laringe}
		La laringe es el órgano de la voz. Actúa como conducto para el paso del aire entre la faringe y la tráquea. Se extiende desde el nivel de \textbf{C3} hasta \textbf{C6}, donde se continúa con la tráquea. Está compuesta por \textbf{cartílago tiroides} y tejido blando; el cartílago tiroides forma la prominencia laríngea o \textit{nuez de Adán}. La \textbf{epiglotis} es una estructura delgada de forma de hoja situada por detrás del techo de la lengua y el hueso hioides, por encima de la entrada laríngea.
		
		\paragraph{Tráquea}
		La tráquea constituye la segunda división del sistema respiratorio propiamente dicho. Es un tubo muscular y fibroso que se extiende desde su unión con la laringe (\textbf{C6}) hasta la \textbf{carina} (\textbf{T4}), punto en el que se divide en los bronquios principales derecho e izquierdo.
		
		Sus paredes contienen aproximadamente \textbf{20 anillos cartilaginosos en forma de C}: completos en su cara anterior e incompletos en su cara posterior, donde una \textbf{membrana fibromusculosa posterior} cierra la luz traqueal. Esta membrana es susceptible de laceración ante una \textbf{mala colocación de tubos endotraqueales}, hecho de relevancia clínica en el contexto radiológico.
		
		\subsubsection*{Faringe}
		
		La faringe constituye una vía común para los sistemas respiratorio y digestivo. Se divide en tres regiones:
		
		\paragraph{Nasofaringe}
		Conecta con las fosas nasales. En esta región desembocan las \textbf{trompas de Eustaquio}, lo que explica la relación entre las \textbf{adenoides hipertróficas} y la \textbf{otitis media}.
		
		\paragraph{Orofaringe}
		Conecta con la cavidad oral. Su límite superior está definido por el \textbf{paladar blando} y la \textbf{úvula}.
		
		\paragraph{Laringofaringe}
		Conecta con la laringe y el esófago, extendiéndose desde la \textbf{epiglotis} hasta el inicio esofágico.
		
		\subsubsection*{Tejido Adenoideo}
		
		\textbf{Desarrollo:} Las adenoides comienzan a ser visibles alrededor de los \textbf{6 meses}, con crecimiento más acelerado entre los \textbf{3 y 6 años}, para posteriormente \textbf{atrofiarse durante la adolescencia}.
		
		\textbf{Ubicación:} Se desarrollan en la base del \textbf{hueso esfenoides}, en la región nasofaríngea.
		
		\textbf{Patología:} La \textbf{hipertrofia adenoidea} causa obstrucción del paso aéreo nasofaríngeo, generando síntomas respiratorios obstructivos y favoreciendo la otitis media por compromiso de las trompas de Eustaquio.
		
		\subsubsection*{Consideraciones Técnicas}
		
		\textbf{Técnica:} Se emplea \textbf{alta tensión} (\textit{high peak voltage}). La región presenta un \textbf{contraste intrínseco propio}: la anatomía regional está compuesta por diversas densidades que no requieren medio de contraste externo.
		
		\textbf{Fundamento del contraste natural:}
		\begin{itemize}
			\item El \textbf{aire} de la vía aérea contrasta directamente con las \textbf{partes blandas} circundantes y el \textbf{cartílago}.
			\item Estructuras radiosensibles adyacentes --- \textbf{glándula tiroides} y \textbf{porciones del esternón} --- deben considerarse en la protección radiológica.
			\item La presencia de un \textbf{cuerpo extraño} se evidencia como una densidad sobreimpuesta sobre la columna de aire.
		\end{itemize}
		
		\clearpage
	
	\twocolumn
		
%%%%%%%%%%%%%%%%%%%%%%%%%%%%%%%%%%%%%%%%%%%%%%%%%%%%%%%%%%%
%  SECCIÓN: Vías Aéreas
%  Archivo: secciones/vias_aereas.tex
%  Sin preámbulo ni \begin{document} — lo maneja main.tex
	%%%%%%%%%%%%%%%%%%%%%%%%%%%%%%%%%%%%%%%%%%%%%%%%%%%%%%%%%%%
	
	\section{Vías Aéreas}
	
	\subsection{Vías Respiratorias Superiores AP
		(Faringe, Laringe y Tráquea)}
	
	\subsubsection*{1. Identificación}
	
	\textbf{Tipo:} Proyección básica de frente. Proyección de referencia para el estudio de las vías aéreas superiores.
	
	\textbf{Indicaciones clínicas principales:}
	\begin{itemize}
		\item Afonía o disfonía
		\item Disfagia (dificultad para tragar)
		\item Cuerpos extraños en vía aérea superior
		\item Desplazamiento traqueal por masas cervicales o mediastínicas
	\end{itemize}
	
	\subsubsection*{2. Técnica Base}
	
	\textbf{Receptor de imagen:} 24 $\times$ 30 cm. Orientación \textbf{longitudinal}, en relación al eje mayor de la región.
	
	\textbf{DFRI:} 100 cm (1 metro). Distancia estándar para esta proyección.
	
	\textbf{Técnica:} Homogeneizada (alta tensión --- \textit{high peak voltage}). Se busca una \textbf{escala de grises larga} que permita visualizar simultáneamente el aire de la vía aérea superior y las partes blandas del cuello. El aumento del kVp favorece la \textbf{radioprotección del paciente} al disminuir la radiación absorbida en la zona tiroidea, que es radiosensible.
	
	\textbf{Bucky:} Sí. Se utiliza Bucky mural dado que la proyección se realiza con el paciente de pie.
	
	\textbf{Consideración sobre movimientos:} Deben reducirse al mínimo. Se recomienda utilizar el \textbf{menor tiempo de exposición posible} (foco grueso) para evitar artefactos de movimiento, especialmente en estructuras laríngeas móviles.
	
	\subsubsection*{3. Posicionamiento}
	
	\paragraph{Paciente}
	De pie (bipedestación) o sentado, con la vía aérea extendida. La posición erguida es fundamental: evita que la columna de aire se ``aplane'' o no se identifique correctamente. El \textbf{plano medio sagital} debe coincidir con la línea media del receptor de imagen.
	
	\paragraph{Región de interés}
	Los hombros deben quedar en el \textbf{mismo plano horizontal}, con el peso distribuido por igual en ambos pies (en bipedestación), para evitar rotaciones del cuello y la cabeza. Se \textbf{extiende la cabeza} lo necesario para evitar que la sombra mandibular oculte la zona laríngea. La base del cráneo y la mandíbula deben quedar por encima de la laringe en la imagen resultante; si persiste superposición, se solicita mayor extensión del mentón.
	
	\begin{figure}[h]
		\centering
		\includegraphics[width=0.7\linewidth]{imagenes/traquea}
		\caption{}
		\label{fig:traquea}
	\end{figure}
	
	
	\paragraph{Rayo Central}
	Perpendicular al plano del receptor. Punto de entrada: \textbf{a la altura de la prominencia laríngea o inmediatamente por debajo de la misma}.
	
	\paragraph{Receptor de imagen}
	Centrado de modo que el campo incluya desde la \textbf{superposición del maxilar inferior y la base del cráneo} hasta los \textbf{vértices pulmonares y el mediastino superior} (laringe y tráquea completas).
	
	\paragraph{Respiración}
	\textbf{Inspiración profunda y lenta.} Llena la cavidad aérea de aire, mejorando el contraste intrínseco entre la columna de aire y las partes blandas circundantes.
	
	\subsubsection*{4. Anatomía Visible}
	
	\textbf{Faringe:} nasofaringe, orofaringe y laringofaringe como columna aérea supraglótica.
	
	\textbf{Laringe:} visible desde \textbf{C3 hasta C6}; se identifica el cartílago tiroides y la epiglotis. Órgano de la voz, compuesto por cartílago tiroides y tejido blando.
	
	\textbf{Tráquea:} visible desde \textbf{C6 hasta T4} (carina). Tubo muscular y fibroso con aproximadamente \textbf{20 anillos cartilaginosos en forma de C}, completos anteriormente e incompletos posteriormente (membrana fibromusculosa posterior). Esta membrana es susceptible de laceración ante mala colocación de tubos endotraqueales.
	
	\textbf{Mediastino superior} y \textbf{vértices pulmonares} incluidos en el margen inferior del campo.
	
	\subsubsection*{5. Criterios de Evaluación}
	
	\textbf{Área incluida:} Desde la superposición del maxilar inferior y la base del cráneo hasta los vértices pulmonares y el mediastino superior.
	
	\textbf{Posicionamiento correcto:} La base del cráneo y la mandíbula deben quedar \textbf{por encima de la laringe}, sin solaparse sobre ella. Si el maxilar inferior se superpone a la laringe, se requiere mayor extensión del mentón.
	
	\textbf{Ausencia de rotación:} Las \textbf{ramas mandibulares deben ser simétricas} y las apófisis espinosas deben aparecer centradas en el cuerpo vertebral. Cuello sin giro lateral.
	
	\textbf{Penetración y contraste:} La densidad radiográfica debe permitir la visualización nítida de las \textbf{estructuras faringolaríngeas}, diferenciando la columna de aire de las partes blandas del cuello.
	
	\subsubsection*{6. Puntos Críticos}
	
	\textbf{Superposición mandibular:} Si la sombra del maxilar inferior oculta la zona laríngea, se corrige solicitando mayor extensión del mentón.
	
	\textbf{Rotación del cuello:} Genera asimetría de las estructuras laríngeas y traqueales. Se verifica con la simetría de las ramas mandibulares y la centración de las apófisis espinosas.
	
	\textbf{Adaptaciones:} En pacientes con limitación para la bipedestación, puede realizarse en decúbito supino, aunque la columna de aire puede no distenderse de forma óptima.
	
	\textbf{Trucos:} Esta proyección es visualmente similar al frente de columna cervical; la diferencia clave radica en la \textbf{técnica} (homogeneizada vs.\ técnica de partes óseas) y en el \textbf{rayo central}: en la proyección de columna cervical el RC se angula cranealmente, mientras que aquí va \textbf{estrictamente perpendicular} al receptor.
	
	\onecolumn
	
	\begin{figure}[h]
		\centering
		\includegraphics[width=0.7\linewidth]{imagenes/traquea2}
		\caption{}
		\label{fig:traquea2}
	\end{figure}
\twocolumn
	
\clearpage
	%----------------------------------------------------------
	%  Para agregar otro enfoque duplicá el bloque
	%  desde \subsection hasta \textbf{Trucos}
	%----------------------------------------------------------

%%%%%%%%%%%%%%%%%%%%%%%%%%%%%%%%%%%%%%%%%%%%%%%%%%%%%%%%%%%
%  SECCIÓN: Vías Aéreas
%  Archivo: secciones/vias_aereas.tex
%  Sin preámbulo ni \begin{document} — lo maneja main.tex
	%%%%%%%%%%%%%%%%%%%%%%%%%%%%%%%%%%%%%%%%%%%%%%%%%%%%%%%%%%%
	
	\section{Vías Aéreas}
	
	\subsection{Proyección Lateral Vías Respiratorias Superiores (Faringe, Laringe y Tráquea)}
	
	\subsubsection*{1. Identificación}
	
	\textbf{Tipo:} Proyección básica de perfil. Complementaria a la proyección AP de vías aéreas superiores.
	
	\textbf{Indicaciones clínicas principales:}
	\begin{itemize}
		\item Descartar laringitis o epiglotitis (especialmente en niños)
		\item Masas o tumores cervicales
		\item Trauma de cuello
		\item Evaluación del diámetro anteroposterior de la vía aérea
		\item Detección de cuerpos extraños
		\item Desplazamientos y compresiones traqueales
	\end{itemize}
	
	\subsubsection*{2. Técnica Base}
	
	\textbf{Receptor de imagen:} 24 $\times$ 30 cm o 30 $\times$ 35 cm. Orientación \textbf{longitudinal} en relación al eje mayor de la región. Borde superior a nivel de la prominencia laríngea.
	
	\textbf{DFRI:} 180 cm recomendado bibliográficamente; se acepta 100 cm. La distancia entre la región anatómica y el receptor genera \textbf{magnificación}, por lo que siempre es preferible aumentar la distancia para mejorar la fidelidad geométrica de la imagen.
	
	\textbf{Técnica:} Homogeneizada. Se busca una \textbf{escala de grises larga} que permita diferenciar el aire de la vía aérea superior de las partes blandas del cuello. El detalle óseo vertebral \textbf{no es requisito} en esta proyección.
	
	\textbf{Bucky:} Sí. Se utiliza Bucky mural.
	
	\subsubsection*{3. Posicionamiento}
	
	\paragraph{Paciente}
	De pie (bipedestación) o sentado, en posición \textbf{lateral estricta}. El \textbf{plano medio sagital} debe quedar paralelo al receptor de imagen y el \textbf{plano coronal} centrado en la línea media del receptor. El peso debe distribuirse equitativamente entre ambos pies.
	
	\paragraph{Región de interés}
	Los \textbf{hombros} deben bajarse al máximo posible, quedando en el \textbf{mismo plano transversal} sin rotación, para evitar su superposición sobre la vía aérea superior. Si el paciente tiene dificultad, se le solicita que tome sus manos por detrás de la espalda para rotar los hombros hacia posterior. El hombro adyacente puede apoyarse sobre el Bucky de pared. Se extiende \textbf{ligeramente la cabeza} del paciente, solicitándole que fije la mirada en un objeto al frente para inmovilizar la región.
	
	\begin{figure}[h]
		\centering
		\includegraphics[width=0.7\linewidth]{imagenes/traquea perfil}
		\caption{}
		\label{fig:traquea-perfil}
	\end{figure}
	
	
	\paragraph{Rayo Central}
	Perpendicular al plano del receptor. Punto de entrada: \textbf{prominencia laríngea (C4--C5)} según referencia clínica, o a nivel de \textbf{C6--C7} entre la escotadura yugular y el plano medio coronal según referencia técnica.
	
	\paragraph{Receptor de imagen}
	Centrado en la \textbf{vía aérea superior}, no en la columna vertebral. El campo debe abarcar desde la \textbf{base del cráneo} hasta la \textbf{parte superior del mediastino} (C1--T1), incluyendo laringe y tráquea completas.
	
	\paragraph{Respiración}
	\textbf{Inspiración profunda y lenta.} Garantiza que la tráquea y la laringe estén llenas de aire, optimizando el contraste intrínseco entre la columna de aire y las partes blandas adyacentes.
	
	\subsubsection*{4. Anatomía Visible}
	
	Laringe en perfil desde \textbf{C3 hasta C6}: cartílago tiroides y tejido blando. Epiglotis especialmente relevante en pediatría (diagnóstico de epiglotitis). Tráquea desde \textbf{C6 hasta T4} (carina): tubo muscular y fibroso con aproximadamente 20 anillos cartilaginosos en forma de C, completos anteriormente e incompletos posteriormente (membrana fibromusculosa posterior susceptible de laceración ante mala colocación de tubos endotraqueales). Espacio retrofaríngeo libre, sin engrosamiento de partes blandas. Luz laríngea visible sin superposición de la columna vertebral.
	
	\subsubsection*{5. Criterios de Evaluación}
	
	\textbf{Posicionamiento correcto:} Campo desde la base del cráneo hasta la parte superior del mediastino (C1--T1). Las diagonales trazadas desde los vértices de la imagen deben converger sobre la \textbf{vía aérea superior}, no sobre la columna cervical. Laringe y tráquea llenas de aire.
	
	\textbf{Ausencia de rotación:} Las \textbf{sombras mandibulares} deben estar superpuestas entre sí; si aparecen desfasadas, indica rotación. Los \textbf{hombros} deben quedar en el mismo plano transversal; si están desfasados, también indica rotación.
	
	\textbf{Penetración y contraste:} La técnica homogeneizada debe permitir la correcta visualización de partes blandas y diferenciar el aire de las estructuras faringolaríngeas. No se exige visualización del detalle óseo vertebral.
	
	\subsubsection*{6. Puntos Críticos}
	
	\textbf{Superposición de hombros:} Oculta el lumen traqueal inferior. Se corrige pidiendo al paciente que baje los hombros o que tome sus manos por detrás de la espalda.
	
	\textbf{Rotación:} La asimetría de las sombras mandibulares o el desfase de los hombros delatan rotación. Ambos deben quedar en sus respectivos planos.
	
	\textbf{Adaptaciones:} En pacientes que no pueden sostenerse de pie, puede realizarse en sedestación manteniendo los mismos criterios de posicionamiento.
	
	\textbf{Trucos:} Ante una imagen de perfil cervical en examen, trazar las diagonales del campo: si convergen sobre la vía aérea y las apófisis espinosas están cortadas, se trata de un perfil de vías aéreas superiores. Si fuera un perfil de columna cervical con apófisis espinosas cortadas, no tendría valor diagnóstico (en columna cervical también debe incluirse el conducto auditivo externo). Diferencias clave con el perfil de columna cervical: en vías aéreas el centrado es sobre la vía aérea (no la columna), el límite superior es más bajo (no requiere el conducto auditivo externo), la técnica es homogeneizada (no para visualizar hueso), las apófisis espinosas pueden cortarse, y el rayo central va \textbf{perpendicular} (no angulado).
	
	\begin{figure*}[h]
		\centering
		\includegraphics[width=0.7\linewidth]{imagenes/traquea perfil2}
		\caption[Flechas blancas: esofago, flechas negras: traquea]{}
		\label{fig:traquea-perfil2}
	\end{figure*}
	

\clearpage
	%----------------------------------------------------------
	%  Para agregar otro enfoque duplicá el bloque
	%  desde \subsection hasta \textbf{Trucos}
	%----------------------------------------------------------

%%%%%%%%%%%%%%%%%%%%%%%%%%%%%%%%%%%%%%%%%%%%%%%%%%%%%%%%%%%
%  SECCIÓN: Vías Aéreas
%  Archivo: secciones/vias_aereas.tex
%  Sin preámbulo ni \begin{document} — lo maneja main.tex
%%%%%%%%%%%%%%%%%%%%%%%%%%%%%%%%%%%%%%%%%%%%%%%%%%%%%%%%%%%


\section{Vías Aéreas}

\subsection{Cavum Perfil}

\subsubsection*{1. Identificación}

\textbf{Tipo:} Proyección básica para estudio pediátrico del cavum nasofaríngeo. \textit{[No se indica proyección complementaria en la información suministrada]}.

\textbf{Indicaciones clínicas principales:}
\begin{itemize}
	\item Hipertrofia adenoidea (vegetaciones)
	\item Obstrucción nasal, rinolalia (voz nasal), ronquidos y apnea del sueño
	\item Otitis media crónica, sinusitis crónica e hipoacusia
	\item Sospecha de neoplasia nasofaríngea
\end{itemize}

\subsubsection*{2. Técnica Base}

\textbf{Receptor de imagen:} 18 $\times$ 24 cm. Orientación \textbf{longitudinal} en paciente pediátrico.

\textbf{DFRI:} entre 100 y 140 cm. Rango que minimiza la magnificación de la nasofaringe y reduce la dosis absorbida en población pediátrica.

\textbf{Técnica:} homogeneizada. Permite visualizar simultáneamente aire, partes blandas y estructuras óseas; kVp alto para penetrar el cráneo combinado con mAs bajos para reducción de dosis en el paciente pediátrico.

\textbf{Factores técnicos:} kVp \textbf{\textit{[dato no suministrado]}} --- mAs \textbf{\textit{[dato no suministrado]}}. \textit{[Método de medición del espesor no suministrado]}.

\textbf{Bucky:} con Bucky y grilla.

\subsubsection*{3. Posicionamiento}

\paragraph{Paciente}
De pie o sentado incorporado; también puede realizarse en decúbito supino con la cara orientada hacia el lado de la mesa. Plano medio sagital paralelo al receptor. Se apoya el lado más cómodo para el paciente; en casos pediátricos puede requerirse colaboración de un familiar para inmovilización.

\paragraph{Región de interés}
Cuello en ligera extensión. Boca cerrada con labios en reposo. Solicitar al paciente que fije la mirada en un punto fijo para evitar rotación e inclinación. El ángulo de la mandíbula no debe superponerse a la columna aérea del cavum. Si se requiere asistencia, un familiar sujeta el cráneo con una mano para evitar rotación y lateralización, y con los dedos de la otra mano sujeta el mentón.

\paragraph{Rayo Central}
Perpendicular al plano del receptor. Punto de entrada: \textbf{2 dedos por delante y 2 dedos por debajo del conducto auditivo externo}.

\paragraph{Receptor de imagen}
Centrado en la región de la nasofaringe/adenoides. Inclusión anatómica: superiormente silla turca y seno esfenoidal (techos orbitarios), la nariz hasta el paladar duro, las tres regiones de la faringe (naso-, oro- y laringofaringe), y hacia posterior desde C1 a C3. No debe incluir hasta T4; el encuadre es más superior que en vías aéreas superiores estándar.

\paragraph{Respiración}
Inspiración nasal tranquila con boca cerrada. La respiración nasal llena la nasofaringe de aire, generando contraste natural entre la columna aérea, el tejido adenoideo y el paladar blando. Se realiza con boca abierta únicamente cuando el médico solicitante lo indica expresamente en el pedido.

\subsubsection*{4. Anatomía Visible}

Paso aéreo continuo desde las fosas nasales hasta la laringofaringe. Tejido adenoideo en la nasofaringe (base del esfenoides). Paladar blando y úvula como referencia de límite entre nasofaringe y orofaringe. Epiglotis delimitando el inicio de la laringofaringe. Silla turca de perfil. Ramas mandibulares superpuestas entre sí. Conducto auditivo externo visible como imagen radiolúcida. Borde cutáneo facial y porción anterior cervical incluidos en el campo.

\subsubsection*{5. Criterios de Evaluación}

\textbf{Posicionamiento correcto:} silla turca de perfil; techos orbitarios sin doble contorno; correcta extensión del mentón con el ángulo mandibular sin superponerse a la columna aérea; región incluida más superior que en vías aéreas estándar, sin llegar a T4.

\textbf{Ausencia de rotación:} ramas mandibulares superpuestas entre sí; conductos auditivos externos superpuestos. Rotación en plano vertical genera doble contorno mandibular de anterior a posterior (PMS no paralelo al receptor). Inclinación en plano sagital genera doble contorno mandibular con una rama superior y otra inferior.

\textbf{Penetración y contraste:} técnica homogeneizada con diferenciación visible entre densidad aérea nasofaríngea, partes blandas (paladar blando, tejido adenoideo) y estructuras óseas (base de cráneo, columna cervical alta).

\subsubsection*{6. Puntos Críticos}

\textbf{Boca abierta sin indicación médica:} el paladar blando se eleva, creando una falsa impresión de obstrucción nasofaríngea que puede conducir a un falso diagnóstico de hipertrofia adenoidea.

\textbf{Rotación del paciente:} la falta de paralelismo entre el PMS y el receptor puede simular hipertrofia adenoidea por superposición asimétrica de tejidos; verificar superposición de ramas mandibulares y CAE antes de exponer.

\textbf{Adaptaciones:} en paciente que no colabora o llora, un familiar puede inmovilizar sujetando el cráneo y el mentón, siempre manteniéndose fuera del campo de radiación y sin interferir con la región de interés radiográfico. En decúbito supino, la cara se orienta lateralmente hacia la mesa.

\textbf{Trucos:} si el paciente traga durante la exposición el paladar blando se eleva, invalidando el estudio; coordinar la toma durante respiración nasal tranquila. Para confirmar que no es un perfil de vías aéreas superiores estándar: el encuadre incluye mayor porción craneal, se identifican estructuras óseas pediátricas y está centrado en la vía aérea superior, sin descender hasta T4.

\clearpage

\onecolumn
\begin{figure}
	\centering
	\includegraphics[width=1\linewidth]{imagenes/cavum}
	\caption{Cavum}
	\label{fig:cavum}
\end{figure}
\twocolumn

\clearpage
%----------------------------------------------------------
%  Para agregar otro enfoque duplicá el bloque
%  desde \subsection hasta \textbf{Trucos}
%----------------------------------------------------------